\chapter{研究结果讨论}
\label{chap:discussion}

\section{方法有效性分析}
\label{sec:discussion:effectiveness}

本节结合第~\ref{chap:experiment}~章实验结果，从数据处理、票房建模、推荐服务与检索统计四个层面，讨论 CineScope 各方法的有效性。

\subsection{数据处理有效性}

ETL 流水线在样本规模与字段质量之间取得了较好平衡。原始电影记录 9742 部，经 TMDb 详情匹配后保留 9727 部，仅丢弃 15 部无法获得有效元数据的影片，保留率达 99.8\%。在此基础上，全部终版电影均具有完整 TMDb 标识，财务元数据成功匹配 9386 部，其中 4 部通过 IMDb 回退策略完成补全；非零票房样本 6235 条、非零预算样本 5844 条，为后续回归建模提供了足够训练规模。

从缺失情况看，简介缺失 35 部（0.36\%）、片长缺失 41 部（0.42\%），关键字段整体可控。评分与标签也随电影主表同步裁剪：终版数据集包含 100811 条评分、3678 条标签，覆盖 610 名用户。该结果说明，基于质量报告的清洗规则能够在不显著牺牲样本量的前提下，为建模与服务层提供稳定、可审计的数据基础。

\subsection{票房预测有效性}

票房预测实验采用 $\log(1+\text{revenue})$ 作为训练目标，在 6235 条非零票房样本中按 8:2 划分训练集（4988）与验证集（1247），共构建 66 维特征。表~\ref{tab:revenue-metrics} 所示结果表明，XGBoost 与 LightGBM 在 $\log$ 空间均取得约 0.64--0.65 的 $R^2$，集成后验证集指标为 $\text{RMSE}_{\log}=1.275$、$\text{MAE}_{\log}=0.857$、$R^2_{\log}=0.652$。

进一步比较单模型与集成效果：XGBoost 单独训练时 $R^2_{\log}=0.652$，LightGBM 为 $0.641$，说明在当前特征体系下 XGBoost 更具优势；按验证集 $\text{RMSE}_{\log}$ 搜索后，集成权重确定为 XGBoost 0.88、LightGBM 0.12，集成模型在 $R^2_{\log}$ 上较 LightGBM 单独提升约 0.011，同时保持与 XGBoost 相当的误差水平。在原始票房空间，集成模型验证集 $\text{RMSE} \approx 9.53 \times 10^7$ 美元、$\text{MAE} \approx 3.94 \times 10^7$ 美元；由于票房分布高度右偏，$\log$ 空间指标更能反映模型对票房数量级的把握能力。

\subsection{推荐系统有效性}

推荐模块离线产物覆盖全部 9727 部电影，协同过滤侧覆盖 610 用户与 100811 条评分；服务健康检查返回 \texttt{artifacts\_ready: true}，说明向量化器、稀疏矩阵与 KNN 模型均已成功加载。在线实验中，以 ``Toy Story'' 为种子调用内容推荐接口，系统能够稳定返回包含电影 ID、相似度分数与来源标签的候选列表，表明 \texttt{server $\rightarrow$ recommender} 调用链可用。

在协同过滤场景下，同时传入 \texttt{user\_id=1} 与种子电影时，系统能够融合 Item KNN 与 User KNN 结果；仅传入用户 ID 时，则返回纯个性化列表。以 ``Toy Story'' 为例，内容推荐 Top-1 为同系列续作 \textit{Toy Story 2}（相似度 0.947），协同推荐 Top-1 同为该续作（分数 0.865），但协同路径还会引入 \textit{Shrek}（0.835）、\textit{The Lion King}（0.768）等用户群体偏好相近的作品。结合 98.3\% 的评分矩阵稀疏度来看，协同过滤能够利用有限行为数据生成候选，而内容推荐则为冷启动与种子影片检索提供了补充路径。两路推荐通过独立微服务部署，也使推荐算法迭代与业务后端聚合实现了有效解耦。

\subsection{检索与统计有效性}

电影检索实验验证了统一搜索与组合过滤策略的可用性。关键词 ``Star'' 可召回 174 部相关电影，涵盖 \textit{Star Wars} 系列及标题含 ``star'' 的其他影片；限定 Sci-Fi 类型时，结果规模与 \texttt{genre\_stats.csv} 中记录的 975 部一致；语言别名映射与分页默认值（25 条/页）均按预期工作，说明后端能够在本地 CSV 数据上支撑交互式查询与结果展示。

统计分析结果与特征工程假设相互印证。表~\ref{tab:correlation} 与表~\ref{tab:feature-importance} 共同表明，预算与热度是票房预测的主导信号，而评分均值的重要性相对有限。在 5230 条预算与票房均非零的样本上，预算与票房的 Pearson 相关系数为 0.735，票房与 TMDb 热度的相关系数为 0.487，而票房与评分均值的相关系数仅为 0.057。XGBoost 特征重要性中，\texttt{budget} 与 \texttt{tmdb\_popularity} 合计占比超过 70\%，与相关性结论一致。同时，高预算并不必然带来高口碑，与可视化实验中“高预算低票房离群点”的观察相符。前端 Atlas 通过 \texttt{/stats/*} 接口在线展示的趋势、散点与相关性视图，与上述定量分析保持一致。

\section{误差来源与局限性}
\label{sec:discussion:limitations}

尽管上述实验支持了方法的有效性，系统仍存在若干不足：

\begin{itemize}
    \item \textbf{样本与特征受限}：非零票房样本 6235 条，且未引入导演、演员、档期等外部特征，限制了预测上限；
    \item \textbf{推荐评估不完整}：尚未系统计算 Precision@K、NDCG@K 等离线指标，协同过滤在稀疏场景下仍存在冷启动问题；
    \item \textbf{RAG 覆盖范围有限}：当前仅索引部分电影画像与摘要文档，语义召回能力仍有提升空间；
\end{itemize}

\section{改进方向}
\label{sec:discussion:future}

结合上述局限，后续工作可从以下方向展开：

\begin{enumerate}
    \item \textbf{完善推荐评估}：补充 Precision@K、NDCG@K 等指标，并引入更系统的离线评测流程；
    \item \textbf{增强语义能力}：采用向量模型替换轻量 RAG，提高问答与推荐解释质量；
    \item \textbf{统一混合推荐}：设计场景化动态加权策略，并提供单一 Hybrid 接口；
\end{enumerate}
